% =====================================================================
% CONCLUSION (paper-wide). Auction + mimic-methodology parts are filled;
% negotiation and deception takeaways are PLACEHOLDERS to fill from those
% sections' final results. No em dashes (en dashes in ranges are fine).
% Note: auction recap here is high-altitude on purpose; the detailed version
% (V, accuracy, the transfer numbers) lives in the auction "Synthesis"
% subsubsection. "compare not train" lesson is stated once, in the methodology paragraph.
% =====================================================================

\section{Conclusion}
We evaluated five frontier language models across three strategic environments under incomplete information: a buyer--seller negotiation, a five-bidder sequential auction, and a five-attribute deception game. A consistent thread runs through all three: each model brings a distinct, stable, and human-legible behavioral signature to the game, yet those signatures reflect a small repertoire of simple heuristics rather than deep adaptive reasoning. % CONFIRM this framing once the negotiation and deception results are final. If either environment shows genuine strategic depth, soften the second clause to "legible per-model signatures" and drop the "simple heuristics" claim.

The auction makes the point most sharply: the LLMs reduce to a one-dimensional fair-share reservation-price strategy that is both \emph{distillable}, in that small networks reproduce their play, and \emph{beatable}, in that trivial, budget-aware, and learned policies all outscore them, non-monotonically in nominal sophistication. Against the real models the ranking broadly transfers, with one exception: the policy trained on the mimics loses its lead to the fixed ones.

% PLACEHOLDER -- negotiation takeaway (1-2 sentences). e.g. which models secure
% the highest buyer- and seller-side utility, how the analytical/Bayesian
% baseline ranks, and the dominant bargaining styles.
In negotiation, \textit{[takeaway to be filled from the negotiation results].}

% PLACEHOLDER -- deception takeaway (1-2 sentences). e.g. how deception rates
% and belief misalignment vary across models, and the role of the truthful
% baseline. (Note: the belief-misalignment metric was in progress; only state
% what the final results support.)
In the deception game, \textit{[takeaway to be filled from the deception results].}

Methodologically, we contribute \emph{behavioral mimics}: small networks cloned from a modest sample of real games that stand in for LLMs at negligible cost, making large-scale strategic evaluation affordable. The explicit caveat, borne out by our transfer experiments, is that they are faithful enough to \emph{compare} fixed policies but not to \emph{train} one for deployment against the real models. Our main limitations follow from scale: the auction results rest on a small number of real-LLM games in a single configuration (high precision but limited scenario diversity), and the real-LLM transfer runs, though now spanning every policy, are only five auctions each. Future work should broaden scenario diversity, run powered real-LLM transfer experiments, and test whether the mimic approach extends to the richer action spaces of the negotiation and deception settings.
